% Source: https://artofproblemsolving.com/community/c3h3186838_2022_abmc_speed_contest__actonboxborough_math_competition
% Stable identity: aops-topic-3186838-q20
Brian is shooting arrows at a target, made of concentric circles of radius $1$, $2$, $3$, and $4$. He gets $10$ points for hitting the innermost circle, $8$ for hitting between the smallest and second smallest circles, $5$ for between the second and third smallest circles, $2$ points for between the third smallest and outermost circle, and no points for missing the target. Assume for each shot he takes, there is a $20\%$ chance Brian will miss the target, but otherwise the chances of hitting each target are proportional to the area of the region. The chance that after three shots, Brian will have scored $15$ points can be expressed as $\dfrac{m}{n}$ for relatively prime positive integers $m, n$. Find $m + n$. Assume the three shots are independent.
